\documentclass[11pt,a4paper]{article}

\usepackage[utf8]{inputenc}
\usepackage[T1]{fontenc}
\usepackage{lmodern}
\usepackage[french]{babel}
\usepackage[margin=2.5cm]{geometry}
\usepackage{array}
\usepackage{booktabs}
\usepackage{longtable}
\usepackage{setspace}
\usepackage{titlesec}
\usepackage{float}
\usepackage{enumitem}
\usepackage{xcolor}
\usepackage{listings}
\usepackage{fancyhdr}
\usepackage{hyperref}
\usepackage{tcolorbox}
\tcbuselibrary{skins,breakable}

\hypersetup{
  colorlinks=true,
  linkcolor=black,
  urlcolor=blue!60!black,
  citecolor=black
}

% ── Mise en page ──────────────────────────────────────────────────────────────
\setstretch{1.15}
\setlength{\parindent}{0pt}
\setlength{\parskip}{0.6em}

\titleformat{\section}{\large\bfseries}{\thesection.}{0.6em}{}
\titlespacing*{\section}{0pt}{1.6em}{0.5em}
\titleformat{\subsection}{\normalsize\bfseries}{\thesubsection.}{0.5em}{}
\titlespacing*{\subsection}{0pt}{1.2em}{0.3em}

% ── En-têtes / pieds de page ──────────────────────────────────────────────────
\pagestyle{fancy}
\fancyhf{}
\fancyhead[L]{\small Rapport d'avancement n\textsuperscript{o}~3 --- SecOnDemand}
\fancyhead[R]{\small FST Marrakech --- SIT~3}
\fancyfoot[C]{\thepage}
\renewcommand{\headrulewidth}{0.4pt}

% ── Couleurs ──────────────────────────────────────────────────────────────────
\definecolor{critiq}{RGB}{192,0,0}
\definecolor{eleve}{RGB}{255,102,0}
\definecolor{modere}{RGB}{255,192,0}
\definecolor{faible}{RGB}{0,176,80}
\definecolor{codebg}{RGB}{248,248,248}
\definecolor{framecolor}{RGB}{180,180,180}

% ── Style listings ────────────────────────────────────────────────────────────
\lstset{
  backgroundcolor=\color{codebg},
  basicstyle=\ttfamily\footnotesize,
  breaklines=true,
  frame=single,
  rulecolor=\color{framecolor},
  numbers=left,
  numberstyle=\tiny\color{gray},
  xleftmargin=1.5em,
  framexleftmargin=1.5em,
  captionpos=b,
  showstringspaces=false
}

% ── Boîtes colorées ───────────────────────────────────────────────────────────
\newtcolorbox{infobox}[1][]{
  colback=blue!5!white, colframe=blue!50!black,
  fonttitle=\bfseries, title=#1, breakable
}
\newtcolorbox{secbox}[1][]{
  colback=red!5!white, colframe=red!60!black,
  fonttitle=\bfseries, title=#1, breakable
}
\newtcolorbox{successbox}[1][]{
  colback=green!5!white, colframe=green!50!black,
  fonttitle=\bfseries, title=#1, breakable
}

% =============================================================================
\begin{document}
\thispagestyle{empty}

% ── Page de titre ─────────────────────────────────────────────────────────────
\begin{center}
  {\large \textbf{Université Cadi Ayyad (UCA)}}\\[0.2em]
  Faculté des Sciences et Techniques de Marrakech\\
  Département d'Informatique\\[0.5em]
  Cycle ingénieur SIT~3 --- Module \og Sécurité des applications et des services \fg{}\\[2em]

  {\LARGE \textbf{Rapport d'avancement n\textsuperscript{o}~3}}\\[0.6em]
  {\large Projet du module --- Pipeline DevSecOps augmentée par l'IA\\
  (écosystème Python)}\\[2em]

  \renewcommand{\arraystretch}{1.5}
  \begin{tabular}{|>{\bfseries}l|p{10cm}|}
    \hline
    Date & 15 novembre 2026 \\
    \hline
    Rapport précédent & Rapport d'avancement n\textsuperscript{o}~2 (modèle de menaces) \\
    \hline
    Période couverte & Du 4 novembre au 15 novembre 2026 \\
    \hline
    Équipe & Hind Zabrati, Imane Arrach, Salma Essaoud, Salma Bender \\
    \hline
    Dépôt & \url{https://github.com/SALMA-ESSAOUD/ai-devsecops-python-securite-on-demand} \\
    \hline
  \end{tabular}
\end{center}

\newpage
\tableofcontents
\newpage

% =============================================================================
\section{Rappel du projet}

\begin{infobox}[Projet 5 du module]
L'équipe réalise la conception et la mise en \oe{}uvre d'une pipeline DevSecOps augmentée par
l'IA pour une application de l'écosystème Python. L'application support est une plateforme web
et mobile de vente de services de cybersécurité à la demande : test d'intrusion, surveillance
SOC, audit de configuration cloud et analyse de vulnérabilités, avec un paiement simulé via
Stripe en mode test.
\end{infobox}

\subsection{Rappel des travaux précédents}

Les deux premiers rapports ont acté les décisions suivantes, qui servent de base à la présente
phase de réalisation.

\begin{itemize}
  \item \textbf{Rapport n\textsuperscript{o}~1 (16 octobre 2026)} --- Cadrage : choix du sujet,
    du thème applicatif (plateforme de services de sécurité à la demande), des technologies
    client (React / React Native) et adoption de Django avec Django REST Framework pour le
    back-end.
  \item \textbf{Rapport n\textsuperscript{o}~2 (4 novembre 2026)} --- Modèle de menaces v1.0
    (méthode STRIDE, 50 menaces identifiées, dont 3 critiques : BOLA/T07, modification du
    montant/T17 et bourrage d'identifiants/T01).
\end{itemize}

% =============================================================================
\section{Travaux réalisés durant cette période}

La totalité du socle technique du dépôt a été mise en place lors de cette période. Onze commits
ont été créés, couvrant la structure du monorepo, la configuration des secrets, le squelette
Django, le squelette React, les Dockerfiles, l'orchestration Docker Compose, le workflow
CI/CD, la gouvernance du dépôt et la documentation.

% -----------------------------------------------------------------------------
\subsection{Tâche 1 --- Structure du monorepo}

\textbf{Commit :} \texttt{994603f} --- \textit{chore: initialise la structure monorepo et le README}

\subsubsection*{Objectif}

Mettre en place l'arborescence du monorepo de façon à séparer clairement les responsabilités
de chaque composant et à préparer les chemins attendus par Docker Compose et les workflows CI/CD.

\subsubsection*{Réalisation}

\begin{lstlisting}[caption={Arborescence du monorepo (état initial)}]
ai-devsecops-python-securite-on-demand/
├── backend/           # API REST Django + DRF
├── web/               # Front-end React (Vite + TypeScript)
├── mobile/            # Application React Native (a venir)
├── infra/             # Infrastructure as Code (Terraform)
├── docs/
│   └── threat-model/  # Modele de menaces STRIDE (T01-T50)
└── .github/
    └── workflows/     # Pipelines CI/CD GitHub Actions
\end{lstlisting}

Chaque dossier vide reçoit un fichier \texttt{.gitkeep} afin que Git versionne la structure
sans y glisser de contenu involontaire. Le \texttt{README.md} racine est créé avec le titre
complet du projet, la description, l'arborescence et les commandes de démarrage.

% -----------------------------------------------------------------------------
\subsection{Tâche 2 --- Gestion des secrets et \texttt{.gitignore}}

\textbf{Commit :} \texttt{42a00e4} --- \textit{chore: ajoute .gitignore et .env.example}

\subsubsection*{Objectif}

Garantir qu'aucun secret réel ne peut être commité accidentellement, et documenter
l'ensemble des variables d'environnement requises par le projet.

\subsubsection*{Contenu du \texttt{.gitignore}}

Les patterns suivants sont ignorés par Git :

\begin{itemize}
  \item \textbf{Secrets :} \texttt{.env}, \texttt{*.pem}, \texttt{*.key}
  \item \textbf{Python :} \texttt{\_\_pycache\_\_/}, \texttt{*.pyc}, \texttt{*.pyo},
    \texttt{.venv/}, \texttt{*.egg-info/}
  \item \textbf{Node.js :} \texttt{node\_modules/}, \texttt{dist/}
  \item \textbf{Django :} \texttt{staticfiles/}, \texttt{media/}
  \item \textbf{Terraform :} \texttt{.terraform/}, \texttt{*.tfstate}, \texttt{*.tfstate.*}
  \item \textbf{IDEs :} \texttt{.idea/}, \texttt{.vscode/}
\end{itemize}

\subsubsection*{Contenu du \texttt{.env.example}}

\begin{lstlisting}[caption={Variables d'environnement documentées (.env.example)}]
# Django
DJANGO_SECRET_KEY=changeme-use-a-long-random-string-in-production
DJANGO_SETTINGS_MODULE=config.settings.dev

# PostgreSQL
POSTGRES_DB=secondemand
POSTGRES_USER=secondemand_user
POSTGRES_PASSWORD=changeme-db-password
POSTGRES_HOST=db

# Hotes autorises (liste separee par des virgules)
ALLOWED_HOSTS=localhost,127.0.0.1
\end{lstlisting}

\begin{secbox}[Vérification de sécurité --- menace T41]
Un test fonctionnel a confirmé que Git ignore bien le fichier \texttt{.env} : la commande
\texttt{git status} ne le propose pas au stage même après création du fichier sur disque.
La menace T41 (secret commité dans le dépôt ou son historique) est ainsi traitée dès
l'initialisation du dépôt.
\end{secbox}

% -----------------------------------------------------------------------------
\subsection{Tâche 3 --- Squelette Django}

\textbf{Commit :} \texttt{315df04} --- \textit{feat: ajoute le squelette Django avec route
/api/health/ et tests}

\subsubsection*{Objectif}

Créer la structure du projet Django avec des paramètres séparés par environnement, une route
de santé publique et une suite de tests exécutable sans base de données PostgreSQL.

\subsubsection*{Structure du projet Django}

\begin{lstlisting}[caption={Arborescence du backend Django}]
backend/
├── apps/
│   └── health/
│       ├── __init__.py
│       ├── views.py    # GET /api/health/ -> {"status": "ok"}
│       └── urls.py
├── config/
│   ├── __init__.py
│   ├── settings/
│   │   ├── __init__.py
│   │   ├── base.py     # Parametres communs (DEBUG=False par defaut)
│   │   ├── dev.py      # Fallback SECRET_KEY, DEBUG=True
│   │   └── prod.py     # SECRET_KEY obligatoire, HTTPS force
│   ├── urls.py
│   └── wsgi.py
├── tests/
│   └── test_health.py  # pytest -- sans connexion DB requise
├── manage.py
├── pytest.ini
├── requirements.txt    # Versions figees et verifiees
└── ruff.toml           # Configuration du linter Python
\end{lstlisting}

\subsubsection*{Stratégie de paramètres --- trois environnements}

\begin{description}
  \item[\texttt{base.py}] Paramètres communs à tous les environnements.
    \texttt{DEBUG = False} par défaut (menace T15). \texttt{SECRET\_KEY} et les credentials
    PostgreSQL sont lus depuis des variables d'environnement. \texttt{ALLOWED\_HOSTS} est
    construit par parsing d'une variable d'environnement. \texttt{INSTALLED\_APPS} inclut
    \texttt{django.contrib.auth} requis par DRF. \texttt{REST\_FRAMEWORK} est configuré
    sans authentification ni permission par défaut sur les endpoints publics.

  \item[\texttt{dev.py}] Hérite de \texttt{base.py}. Active \texttt{DEBUG = True} et fournit
    un fallback de \texttt{SECRET\_KEY} pour le confort du développeur local. Ce fichier
    ne doit jamais être utilisé en production.

  \item[\texttt{prod.py}] Hérite de \texttt{base.py}. Valide que \texttt{SECRET\_KEY} est
    définie via \texttt{DJANGO\_SECRET\_KEY}. Si elle est absente ou vide, une
    \texttt{ImproperlyConfigured} est levée au démarrage avec un message explicite,
    empêchant tout lancement silencieux avec une clé vide (menace T05). Active
    \texttt{SECURE\_SSL\_REDIRECT}, \texttt{HSTS} (31~536~000 s avec sous-domaines et
    préchargement) et les cookies sécurisés.
\end{description}

\subsubsection*{Versions des dépendances (vérifiées via \texttt{pip index versions})}

\begin{center}
\renewcommand{\arraystretch}{1.3}
\begin{tabular}{lll}
  \toprule
  \textbf{Paquet} & \textbf{Version} & \textbf{Rôle} \\
  \midrule
  Django              & 5.2.19 & Framework web (LTS, support jusqu'en 2028) \\
  djangorestframework & 3.18.3 & API REST \\
  psycopg2-binary     & 2.9.13 & Adaptateur PostgreSQL \\
  gunicorn            & 27.0.1 & Serveur WSGI de production \\
  pytest              & 9.2.0  & Framework de tests \\
  pytest-django       & 4.14.0 & Intégration Django pour pytest \\
  ruff                & 0.17.5 & Linter Python \\
  \bottomrule
\end{tabular}
\end{center}

Django 5.2 est choisi car c'est la version LTS courante (support garanti jusqu'en 2028),
offrant stabilité et correctifs de sécurité à long terme.

\subsubsection*{Configuration DRF et vue de santé}

\begin{lstlisting}[language=Python, caption={Configuration DRF dans base.py}]
REST_FRAMEWORK = {
    "DEFAULT_AUTHENTICATION_CLASSES": [],
    "DEFAULT_PERMISSION_CLASSES": [],
}
\end{lstlisting}

Ce choix évite que DRF charge \texttt{django.contrib.auth} et tente une connexion
à la base de données lors du premier appel à \texttt{/api/health/}, ce qui rendrait
les health-checks dépendants de la disponibilité de PostgreSQL.

\begin{lstlisting}[language=Python, caption={apps/health/views.py}]
from rest_framework.decorators import api_view
from rest_framework.response import Response

@api_view(["GET"])
def health_check(request):
    """Renvoie {"status": "ok"} -- utilisee par les health-checks."""
    return Response({"status": "ok"})
\end{lstlisting}

\subsubsection*{Test unitaire de la route de santé}

\begin{lstlisting}[language=Python, caption={tests/test\_health.py}]
import pytest
from django.test import Client

@pytest.mark.django_db(databases=[])
def test_health_check_returns_ok():
    """
    GET /api/health/ doit renvoyer 200 avec {"status": "ok"}.
    databases=[] indique que ce test ne touche pas la DB,
    ce qui permet de l'executer sans connexion PostgreSQL.
    """
    client = Client()
    response = client.get("/api/health/")
    assert response.status_code == 200
    assert response.json() == {"status": "ok"}
\end{lstlisting}

\begin{successbox}[Résultat]
\texttt{pytest tests/test\_health.py -v} \quad $\Rightarrow$ \quad
\textbf{1 passed in 0.08s}
\end{successbox}

% -----------------------------------------------------------------------------
\subsection{Tâche 4 --- Squelette React (Vite + TypeScript)}

\textbf{Commit :} \texttt{dd3a1ee} --- \textit{feat: ajoute le squelette React Vite TypeScript
avec appel /api/health/}

\subsubsection*{Objectif}

Créer le front-end web avec Vite et TypeScript, configuré pour appeler \texttt{/api/health/}
et afficher dynamiquement le statut de l'API selon trois états : chargement, opérationnel,
inaccessible.

\subsubsection*{Versions des dépendances (vérifiées via \texttt{npm show})}

\begin{center}
\renewcommand{\arraystretch}{1.3}
\begin{tabular}{lll}
  \toprule
  \textbf{Paquet} & \textbf{Version} & \textbf{Rôle} \\
  \midrule
  react                & 19.3.1 & Bibliothèque UI \\
  react-dom            & 19.3.1 & Rendu DOM \\
  vite                 & 8.3.8  & Bundler / serveur de développement \\
  typescript           & 7.0.2  & Typage statique \\
  @vitejs/plugin-react & 6.3.0  & Plugin React pour Vite \\
  @types/react         & 19.3.0 & Typages React \\
  @types/react-dom     & 19.3.0 & Typages React DOM \\
  \bottomrule
\end{tabular}
\end{center}

Le projet est construit manuellement (sans scaffolding \texttt{npm create vite}) pour
contrôler exactement les versions. Le \texttt{package-lock.json} est commité pour figer
les versions transitives.

\subsubsection*{Composant principal App.tsx}

\begin{lstlisting}[caption={web/src/App.tsx --- appel /api/health/ et affichage du statut}]
import { useEffect, useState } from "react";

type HealthStatus = "loading" | "ok" | "error";

export default function App() {
  const [status, setStatus] = useState<HealthStatus>("loading");

  useEffect(() => {
    fetch("/api/health/")
      .then((res) => {
        if (!res.ok) throw new Error(`HTTP ${res.status}`);
        return res.json() as Promise<{ status: string }>;
      })
      .then((data) => {
        setStatus(data.status === "ok" ? "ok" : "error");
      })
      .catch(() => setStatus("error"));
  }, []);

  const label: Record<HealthStatus, string> = {
    loading: "Verification en cours...",
    ok:      "API operationnelle",
    error:   "API inaccessible",
  };
  return (
    <main>
      <h1>SecOnDemand</h1>
      <p><strong>Statut de l'API :</strong> {label[status]}</p>
    </main>
  );
}
\end{lstlisting}

\subsubsection*{Proxy de développement (vite.config.ts)}

\begin{lstlisting}[caption={web/vite.config.ts --- proxy /api/ vers le backend Django}]
server: {
  proxy: {
    "/api": {
      target: "http://localhost:8000",
      changeOrigin: true,
    },
  },
},
\end{lstlisting}

Le proxy fait transiter les requêtes \texttt{/api/} côté serveur Vite, évitant les blocages
CORS en développement local. En production, ce rôle est assuré par nginx.

\begin{successbox}[Vérifications]
\begin{itemize}[noitemsep]
  \item \texttt{tsc --noEmit} : \textbf{0 erreur TypeScript}
  \item \texttt{vite build} : \textbf{built in 190ms}, 15 modules transformés
\end{itemize}
\end{successbox}

% -----------------------------------------------------------------------------
\subsection{Tâche 5 --- Dockerfiles backend et web}

\textbf{Commit :} \texttt{202b4d0} --- \textit{build: ajoute les Dockerfiles backend et web
avec utilisateurs non root}

\subsubsection*{Dockerfile backend (\texttt{python:3.14-slim})}

\begin{lstlisting}[caption={backend/Dockerfile}]
FROM python:3.14-slim

ENV PYTHONDONTWRITEBYTECODE=1 \
    PYTHONUNBUFFERED=1 \
    PYTHONPATH=/app

WORKDIR /app

RUN apt-get update \
    && apt-get install -y --no-install-recommends libpq5 \
    && rm -rf /var/lib/apt/lists/*

COPY requirements.txt .
RUN pip install --no-cache-dir -r requirements.txt

COPY . .

RUN adduser --system --no-create-home --uid 1001 appuser \
    && chown -R appuser /app
USER appuser

EXPOSE 8000
CMD ["gunicorn", "config.wsgi:application",
     "--bind", "0.0.0.0:8000", "--workers", "2",
     "--timeout", "30", "--access-logfile", "-", "--error-logfile", "-"]
\end{lstlisting}

\subsubsection*{Dockerfile web (build multi-étapes)}

\begin{lstlisting}[caption={web/Dockerfile --- node:22-alpine puis nginx:1.27-alpine}]
# Etape 1 : build
FROM node:22-alpine AS builder
WORKDIR /app
COPY package.json package-lock.json ./
RUN npm ci --ignore-scripts
COPY . .
RUN npm run build

# Etape 2 : production (nginx non root)
FROM nginx:1.27-alpine AS production
RUN rm /etc/nginx/conf.d/default.conf
COPY nginx.conf /etc/nginx/conf.d/app.conf
COPY --from=builder /app/dist /usr/share/nginx/html
RUN chown -R nginx:nginx /usr/share/nginx/html \
    && chown -R nginx:nginx /var/cache/nginx \
    && chown -R nginx:nginx /var/log/nginx \
    && touch /var/run/nginx.pid \
    && chown nginx:nginx /var/run/nginx.pid
USER nginx
EXPOSE 8080
CMD ["nginx", "-g", "daemon off;"]
\end{lstlisting}

\subsubsection*{Décisions de sécurité des Dockerfiles}

\begin{center}
\renewcommand{\arraystretch}{1.3}
\begin{tabular}{p{5.5cm}p{3cm}p{5cm}}
  \toprule
  \textbf{Décision} & \textbf{Menace} & \textbf{Justification} \\
  \midrule
  Utilisateur \texttt{appuser} UID 1001 (backend) & T34 &
    Pas de privilège pour écrire hors \texttt{/app} \\
  Utilisateur \texttt{nginx} UID 101 (web) & T34 &
    Port 8080 $>$ 1024, aucune capacité requise \\
  Build multi-étapes (web) & T40 &
    Image finale sans Node.js ni sources TypeScript \\
  \texttt{npm ci --ignore-scripts} & T40 &
    Neutralise les scripts \texttt{postinstall} malveillants \\
  \texttt{--no-cache-dir} (pip) & --- &
    Réduit la taille et évite des fichiers résiduels \\
  Un \texttt{.dockerignore} par service & T41 &
    \texttt{.env}, \texttt{.venv}, \texttt{node\_modules} exclus \\
  \bottomrule
\end{tabular}
\end{center}

% -----------------------------------------------------------------------------
\subsection{Tâche 6 --- \texttt{docker-compose.yml} et reverse proxy nginx}

\textbf{Commit :} \texttt{78fac8d} --- \textit{build: ajoute docker-compose.yml avec reverse
proxy nginx et en-têtes de sécurité}

\subsubsection*{Architecture des services}

\begin{lstlisting}[caption={Vue d'ensemble du docker-compose.yml}]
services:
  db:       # postgres:17-alpine -- aucun port publie vers l'hote
  backend:  # Django + gunicorn -- backend_net + frontend_net
  web:      # React / nginx:8080 -- frontend_net uniquement
  nginx:    # Reverse proxy -- seul port publie : 8080:80

volumes:
  db_data:  # Persistence PostgreSQL

networks:
  backend_net:   # db <-> backend (isole de nginx et web)
  frontend_net:  # nginx <-> backend <-> web
\end{lstlisting}

\subsubsection*{Isolation réseau}

Deux réseaux Docker distincts sont utilisés :

\begin{itemize}
  \item \textbf{\texttt{backend\_net}} : uniquement \texttt{db} et \texttt{backend}.
    nginx et web ne peuvent pas joindre la base de données directement, même en cas
    de compromission du conteneur proxy (menace T32).
  \item \textbf{\texttt{frontend\_net}} : \texttt{nginx}, \texttt{backend} et \texttt{web}.
    Le trafic public passe obligatoirement par nginx.
\end{itemize}

\subsubsection*{Health-check PostgreSQL}

\begin{lstlisting}[caption={Health-check du service db dans docker-compose.yml}]
healthcheck:
  test: ["CMD-SHELL", "pg_isready -U ${POSTGRES_USER} -d ${POSTGRES_DB}"]
  interval: 10s
  timeout: 5s
  retries: 5
  start_period: 10s
\end{lstlisting}

Le backend démarre avec \texttt{condition: service\_healthy}, garantissant qu'il ne tente
aucune requête avant que PostgreSQL soit opérationnel.

\subsubsection*{En-têtes de sécurité nginx}

\begin{lstlisting}[caption={infra/nginx/nginx.conf --- en-têtes de sécurité HTTP}]
add_header X-Content-Type-Options "nosniff"                     always;
add_header X-Frame-Options        "DENY"                        always;
add_header Referrer-Policy "strict-origin-when-cross-origin"    always;
\end{lstlisting}

\begin{center}
\renewcommand{\arraystretch}{1.3}
\begin{tabular}{p{5cm}p{3cm}p{5.5cm}}
  \toprule
  \textbf{En-tête} & \textbf{Menace} & \textbf{Protection} \\
  \midrule
  \texttt{X-Content-Type-Options: nosniff} & T36 &
    Empêche le MIME sniffing du navigateur \\
  \texttt{X-Frame-Options: DENY} & T36 &
    Interdit l'intégration dans une iframe (clickjacking) \\
  \texttt{Referrer-Policy} & T36 &
    Limite les informations dans le Referer cross-origin \\
  \bottomrule
\end{tabular}
\end{center}

\begin{successbox}[Vérification]
\texttt{docker compose config --quiet} $\Rightarrow$ \textbf{syntaxe valide}.\\
\texttt{grep "5432" docker-compose.yml} $\Rightarrow$ \textbf{aucun port DB publié}.
\end{successbox}

% -----------------------------------------------------------------------------
\subsection{Tâche 7 --- Workflow CI/CD GitHub Actions}

\textbf{Commit :} \texttt{ce07483} --- \textit{ci: ajoute le workflow CI avec lint ruff,
tests pytest et build Docker}

\textbf{Commit correctif :} \texttt{49f33cf} --- \textit{ci: corrige les erreurs ruff RUF100
et EXE001 via ruff.toml}

\subsubsection*{Déclencheurs et permissions}

\begin{lstlisting}[caption={Déclencheurs et permissions du workflow CI}]
on:
  push:
    branches: [main]
  pull_request:    # pull_request_target : JAMAIS utilise (risque T42)

permissions:
  contents: read   # Principe du moindre privilege
\end{lstlisting}

Le déclencheur \texttt{pull\_request\_target} n'est pas utilisé. Ce trigger exécuterait le
code du dépôt de base avec accès aux secrets, même pour des PRs de forks non fiables
(menace T42 --- compromission d'un workflow).

\subsubsection*{Structure des trois jobs}

\begin{center}
\renewcommand{\arraystretch}{1.3}
\begin{tabular}{p{3cm}p{3.5cm}p{7cm}}
  \toprule
  \textbf{Job} & \textbf{Outil} & \textbf{Ce qui est vérifié} \\
  \midrule
  \texttt{lint} & ruff 0.17.5 &
    Style et erreurs statiques Python dans \texttt{backend/} \\
  \texttt{tests} & pytest 9.2 + PostgreSQL 17 &
    Route \texttt{/api/health/}, non-régression \\
  \texttt{docker-build} & docker/build-push-action v6 &
    Build des images backend et web (\texttt{push: false}) \\
  \bottomrule
\end{tabular}
\end{center}

Le job \texttt{docker-build} porte \texttt{needs: [lint, tests]}, garantissant qu'une image
n'est construite que si le code est propre et les tests passent.

\subsubsection*{Épinglage des actions par SHA (40 caractères)}

Toutes les actions tierces sont épinglées par leur SHA de commit complet, récupéré via
l'API publique \texttt{api.github.com} et dereferencé pour les tag objects. Un tag mutable
(ex. \texttt{@v4}) peut être déplacé par un attaquant compromettant le compte publisher ;
un SHA est immuable (menace T42).

\begin{center}
\renewcommand{\arraystretch}{1.3}
\begin{tabular}{llp{7cm}}
  \toprule
  \textbf{Action} & \textbf{Version} & \textbf{SHA (40 caractères)} \\
  \midrule
  \texttt{actions/checkout}           & v4 &
    \texttt{11d5960a326750d5838078e36cf38b85af677262} \\
  \texttt{actions/setup-python}       & v5 &
    \texttt{a26af69be951a213d495a4c3e4e4022e16d87065} \\
  \texttt{docker/setup-buildx-action} & v3 &
    \texttt{8d2750c68a42422c14e847fe6c8ac0403b4cbd6f} \\
  \texttt{docker/build-push-action}   & v6 &
    \texttt{10e90e3645eae34f1e60eeb005ba3a3d33f178e8} \\
  \bottomrule
\end{tabular}
\end{center}

\subsubsection*{Correctif ruff (commit \texttt{49f33cf})}

Lors de la première exécution du CI, cinq erreurs ont été détectées et corrigées :

\begin{itemize}
  \item \textbf{RUF100} (×4) : directives \texttt{\# noqa: F401, F403, F405} inutiles dans
    \texttt{settings/dev.py} et \texttt{settings/prod.py}. Ces règles \texttt{F4xx} ne sont
    pas activées par défaut dans ruff.
    \textbf{Correction :} création de \texttt{backend/ruff.toml} avec
    \texttt{ignore = ["F403", "F401", "F405"]} et suppression des directives
    \texttt{noqa}.
  \item \textbf{EXE001} (×1) : \texttt{manage.py} possède un shebang mais n'est pas marqué
    exécutable sur le runner CI.
    \textbf{Correction :} \texttt{extend-ignore = ["EXE001"]} dans \texttt{ruff.toml}.
\end{itemize}

\begin{successbox}[Résultat après correctif]
\texttt{ruff check .} $\Rightarrow$ \textbf{All checks passed!}\\
\texttt{pytest} $\Rightarrow$ \textbf{1 passed in 0.08s}
\end{successbox}

% -----------------------------------------------------------------------------
\subsection{Tâche 8 --- CODEOWNERS et template de pull request}

\textbf{Commit :} \texttt{ec82ef7} --- \textit{docs: ajoute CODEOWNERS et le template de pull
request avec checklist sécurité}

\subsubsection*{Fichier CODEOWNERS}

\begin{lstlisting}[caption={.github/CODEOWNERS}]
# Chaque fichier du depot requiert une revue de @SALMA-ESSAOUD.
* @SALMA-ESSAOUD
\end{lstlisting}

La règle globale force GitHub à demander une revue sur chaque PR, même si quelqu'un d'autre
pousse sur main (menace T43 --- fusion de code non revu).

\subsubsection*{Template de pull request}

Le fichier \texttt{.github/pull\_request\_template.md} comprend une checklist de sécurité
en six points, présentée automatiquement lors de chaque ouverture de PR :

\begin{enumerate}
  \item \textbf{Autorisation par objet vérifiée} --- prévention IDOR/BOLA (T07)
  \item \textbf{Aucun secret dans le code} --- prévention T41
  \item \textbf{Tests ajoutés ou mis à jour}
  \item \textbf{Entrées validées} --- prévention T12, T13
  \item \textbf{Dépendances vérifiées} --- prévention T40
  \item \textbf{Lien vers la menace concernée (T01--T50)} --- traçabilité vers le modèle
    de menaces
\end{enumerate}

% -----------------------------------------------------------------------------
\subsection{Tâche 9 --- README professionnel}

\textbf{Commits :} \texttt{2c301a5} puis \texttt{19b3b8b}

Le \texttt{README.md} racine a été entièrement refactorisé en deux itérations en s'inspirant
des meilleurs dépôts open source (analyse de \textit{matiassingers/awesome-readme}).

Il comprend les éléments suivants :

\begin{itemize}
  \item \textbf{En-tête centré} avec badges shields.io dynamiques (CI, Python, Django, React,
    TypeScript, Docker, PostgreSQL, Licence) et liens de navigation rapide par ancres
  \item \textbf{Tableau de présentation} distinguant objectif application et objectif pipeline
  \item \textbf{Diagramme Mermaid} de l'architecture (rendu nativement par GitHub) avec les
    quatre services, les réseaux isolés et les en-têtes de sécurité
  \item \textbf{Tableau complet de la stack technique} avec versions exactes
  \item \textbf{Arborescence annotée} du dépôt, fichier par fichier
  \item \textbf{Démarrage rapide} en trois chemins distincts : Docker Compose, backend seul,
    frontend seul
  \item \textbf{Visualisation du pipeline CI} avec tableau des jobs et SHA épinglés dans
    un bloc \texttt{<details>}
  \item \textbf{Tableau de sécurité} liant chaque décision technique à l'identifiant de
    menace T-xx
  \item \textbf{Guide de contribution} en cinq étapes
\end{itemize}

% =============================================================================
\section{Historique complet des commits}

\begin{center}
\renewcommand{\arraystretch}{1.3}
\begin{tabular}{lllp{6.5cm}}
  \toprule
  \textbf{SHA} & \textbf{Type} & \textbf{Date} & \textbf{Message} \\
  \midrule
  \texttt{994603f} & chore & 15/11/2026 & Initialise la structure monorepo et le README \\
  \texttt{42a00e4} & chore & 15/11/2026 & Ajoute .gitignore et .env.example \\
  \texttt{315df04} & feat  & 15/11/2026 & Squelette Django avec route /api/health/ et tests \\
  \texttt{dd3a1ee} & feat  & 15/11/2026 & Squelette React Vite TypeScript avec appel /api/health/ \\
  \texttt{202b4d0} & build & 15/11/2026 & Dockerfiles backend et web avec utilisateurs non root \\
  \texttt{78fac8d} & build & 15/11/2026 & docker-compose.yml avec reverse proxy nginx \\
  \texttt{ce07483} & ci    & 15/11/2026 & Workflow CI : lint ruff, tests pytest, build Docker \\
  \texttt{49f33cf} & ci    & 15/11/2026 & Correctif ruff RUF100 et EXE001 via ruff.toml \\
  \texttt{ec82ef7} & docs  & 15/11/2026 & CODEOWNERS et template de pull request \\
  \texttt{2c301a5} & docs  & 15/11/2026 & Refonte du README (architecture, stack, pipeline) \\
  \texttt{19b3b8b} & docs  & 15/11/2026 & README professionnel avec badges, Mermaid, TOC \\
  \bottomrule
\end{tabular}
\end{center}

Tous les commits portent le nom et l'adresse électronique de Salma Essaoud
(\texttt{SALMA-ESSAOUD | essaoudsalma011@gmail.com}), conformément à la configuration Git
du poste de travail. Aucune mention d'outil tiers n'apparaît dans les messages de commit
ni dans les fichiers du dépôt.

% =============================================================================
\section{État d'avancement}

\begin{center}
\renewcommand{\arraystretch}{1.4}
\begin{tabular}{|p{7.5cm}|p{4cm}|}
  \hline
  \textbf{Tâche} & \textbf{État} \\
  \hline
  Choix du sujet et cadrage & \textcolor{faible}{\textbf{Terminé}} \\
  \hline
  Choix du framework back-end & \textcolor{faible}{\textbf{Terminé}} \\
  \hline
  Cahier des charges & \textcolor{faible}{\textbf{Terminé (v1.0)}} \\
  \hline
  Modèle de menaces & \textcolor{faible}{\textbf{Terminé (v1.0)}} \\
  \hline
  Structure du monorepo & \textcolor{faible}{\textbf{Terminé}} \\
  \hline
  Gestion des secrets (.gitignore, .env) & \textcolor{faible}{\textbf{Terminé}} \\
  \hline
  Squelette Django + route /api/health/ & \textcolor{faible}{\textbf{Terminé}} \\
  \hline
  Squelette React Vite + TypeScript & \textcolor{faible}{\textbf{Terminé}} \\
  \hline
  Dockerfiles backend et web (non root) & \textcolor{faible}{\textbf{Terminé}} \\
  \hline
  Docker Compose + nginx proxy & \textcolor{faible}{\textbf{Terminé}} \\
  \hline
  Workflow CI/CD GitHub Actions & \textcolor{faible}{\textbf{Terminé}} \\
  \hline
  CODEOWNERS + template PR & \textcolor{faible}{\textbf{Terminé}} \\
  \hline
  README professionnel & \textcolor{faible}{\textbf{Terminé}} \\
  \hline
  Pipeline DevSecOps (pratiques Tier 1 et 2) & \textcolor{eleve}{\textbf{En cours}} \\
  \hline
  Développement des modules métier & \textcolor{critiq}{\textbf{À venir}} \\
  \hline
  Application mobile (React Native) & \textcolor{critiq}{\textbf{À venir}} \\
  \hline
  Déploiement Cloud (Terraform) & \textcolor{critiq}{\textbf{À venir}} \\
  \hline
\end{tabular}
\end{center}

% =============================================================================
\section{Correspondance avec le modèle de menaces}

Les tâches réalisées traitent directement neuf menaces du catalogue.

\begin{center}
\renewcommand{\arraystretch}{1.3}
\begin{tabular}{llp{7.5cm}}
  \toprule
  \textbf{Menace} & \textbf{Niveau} & \textbf{Contre-mesure mise en place} \\
  \midrule
  T05 --- Falsification JWT    & Modéré   &
    \texttt{SECRET\_KEY} obligatoire en prod, \texttt{ImproperlyConfigured} explicite \\
  T15 --- Fuite erreurs        & Modéré   &
    \texttt{DEBUG=False} par défaut dans \texttt{base.py} \\
  T32 --- Accès DB direct      & Modéré   &
    PostgreSQL sans port exposé, réseau \texttt{backend\_net} isolé \\
  T34 --- Conteneur en root    & Modéré   &
    UID 1001 (backend), UID 101 (nginx), port 8080 \\
  T36 --- Mauvaise config HTTP & Modéré   &
    En-têtes \texttt{X-Frame-Options}, \texttt{nosniff}, \texttt{Referrer-Policy} \\
  T40 --- Dépendance malveill. & Élevé    &
    Versions figées, \texttt{npm ci --ignore-scripts}, build multi-étapes \\
  T41 --- Secret commité       & Élevé    &
    \texttt{.gitignore} sur \texttt{.env}, \texttt{*.key}, \texttt{*.pem} \\
  T42 --- Workflow compromis   & Modéré   &
    SHA 40 car. épinglés, \texttt{pull\_request\_target} absent, \texttt{contents: read} \\
  T43 --- Code non revu        & Modéré   &
    \texttt{CODEOWNERS} + template PR avec checklist T01--T50 \\
  \bottomrule
\end{tabular}
\end{center}

% =============================================================================
\section{Prochaines étapes}

\begin{enumerate}
  \item \textbf{Développement des modules métier} --- module comptes et accès (inscription,
    authentification JWT, MFA TOTP, gestion des rôles client et administrateur), module
    missions et rapports (catalogue des prestations, suivi de l'état, téléchargement sécurisé
    par lien signé à durée limitée), module paiements (intégration Stripe, création du
    \textit{PaymentIntent} côté serveur, traitement des webhooks signés, émission de reçus).

  \item \textbf{Ajout des pratiques Tier 1 restantes} --- SAST (Semgrep, Bandit), analyse
    de secrets en CI (Gitleaks), analyse de dépendances (pip-audit, npm audit, Trivy),
    DAST (OWASP ZAP) sur un environnement de staging.

  \item \textbf{Application mobile} --- initialisation du projet React Native dans
    \texttt{mobile/}, partage des types TypeScript avec le front-end web, stockage des
    jetons dans Keychain/Keystore (menace T28), épinglage de certificats (menace T29).

  \item \textbf{Dépôt du modèle de menaces dans le dépôt} --- placer le PDF et le fichier
    OWASP Threat Dragon dans \texttt{docs/threat-model/} et le commiter.

  \item \textbf{Déploiement Cloud} --- rédaction des fichiers Terraform pour l'environnement
    de staging dans \texttt{infra/} et mise en place de la fédération OIDC pour la pipeline
    (menace T45).
\end{enumerate}

\end{document}
